\section{ML for seed selection and elimination ordering}
\label{sec:ml-roles}
To study learned choices in IBP reduction, we use two System 1
models: one for seed selection and one for elimination ordering.
Here a System 1 model is a small machine learning model that learns
from recorded reductions and assigns scores to choices available at
the current step. Training requires targets with known master
expansions, the seeds that supply their equations and the steps used
to reduce them. The implemented models and their completed training
appear in Sections~\ref{sec:seed-model} and
\ref{sec:reduction-model}, followed by their held-out tests.

In a Kira integration, the seed model would act before equation
generation and the elimination model would act when the solver
selects a pivot. Kira would perform the selected algebraic operations.
The present experiments test seed-reference agreement and reduction
of a target coefficient row against a fixed source-equation pool.
Applying the models to Kira's full evolving system remains a separate
step.

\subsection{Selecting seeds before equation generation}
The first model receives the family definition, the target terms and
their index vectors $\boldsymbol a$, the chosen master basis and a
candidate seed list within stated sectors and index bounds. It
assigns each candidate a separate probability of belonging to the
target's recorded sufficient seed set, which allows several seeds
to receive label one.

During training, we pair each target with the seeds used by its
checked reference reduction. Different reductions can use different
sufficient sets, and the stored labels identify the selected reference
set. They do not establish that every omitted seed is unnecessary.
For use inside Kira, a reduction would generate equations from the
predicted seeds and check whether they suffice in the chosen master
basis. An unresolved target would require further seeds. The
seed-model test below measures recovery of the recorded sets; it
does not perform this complete solver test.

\subsection{Choosing elimination steps within the solver}
The second model receives the current coefficient matrix and the
supplied allowed row operations. It assigns a probability to each
operation, selects one operation and reports the selected
probability as its confidence. An exact row update then applies
Eq.~\eqref{eq:row-update}. The next state supplies the information
for the next choice.

Each training example specifies the matrix before the choice, the
allowed pivots, and the recorded source row and integral to eliminate.
The retained later steps permit checks of the complete path and its
arithmetic count, while predictive features use the current state
and supplied actions. In the present data, the source identities stay
fixed and one target row changes during reduction. Extending this
procedure to Kira requires allowed choices that follow Kira's
integral ordering and evolving equations.

The computational benefit depends on the total work: preparing
features and evaluating the model must cost less than the algebraic
work that its choices save. The tests therefore record complete-path
arithmetic and timing as well as next-step agreement. Kira's existing
finite-field reconstruction would remain responsible for obtaining
symbolic coefficients in a later solver integration.

\subsection{Integral data required for training and testing}
To train and test both models, we need family denominators, kinematic
parameters, targets and a master basis together with the exact
coefficients in Eq.~\eqref{eq:reduction}. For a target containing
several integrals, these coefficients specify the master expansion
of the complete expression. Each example also identifies its source
equations, actual IBP seeds and derivative operators. Algebraic
identities, including mass partial fractions, remain separately
identified because they carry no IBP seed.

We construct targets from known master combinations and retain the
operations used to form each target and reduce it back to the same
basis. The elimination records give the coefficients before and after
each step, all allowed source rows, the selected pivot and the
operation counts. These records connect seed selection to the
equations available for elimination.

The known master coefficients provide a reference for checking
complete reduction in the chosen basis, as described in
Section~\ref{sec:ibp-reduction}~\cite{Kira3}. We keep these answers,
the selected labels and future construction histories separate from
the information used to predict a choice.
Section~\ref{sec:integrals-generation} explains how we obtain the
bases, construct the targets and check their seed and elimination
records.
